\section{From finite factors to arbitrary homomorphisms}\label{sec:global}

We now remove the restrictions on the algebra and its representation, and
then pass from inner commutators to arbitrary derivations.  Throughout this
section set
\begin{equation}\label{eq:global-constants}
 C_{\mathrm{fac}}
   =\max\{3c_1+2,\,1+2c_{\mathrm{row}},\,2\},
 \qquad C=3C_{\mathrm{fac}}+2.
\end{equation}
These constants depend only on the numerical constants of the preceding
sections.

\subsection{Representations of factors}

\begin{proposition}\label{prop:factor-representations}
Let $P\subset\B(K)$ be a unital factor acting on a separable Hilbert space.
For every $Y\in\B(K)$, every $h\ge1$, and every $X\in M_h(P)$,
\[
 \norm{[Y^{(h)},X]}
   \le C_{\mathrm{fac}}\,g_P(Y)\,\norm{X}.
\]
\end{proposition}

\begin{proof}
We first treat a type $\mathrm{II}_1$ factor $M$ in an arbitrary normal
representation on a separable Hilbert space.  Its predual is separable.
Every cyclic normal representation embeds into the tracial standard
representation.  Indeed, if $\xi$ is a cyclic vector, its normal positive
functional has a trace density $a\in L^1(M)_+$
\citep[Corollary~5.14 and Remark~5.15]{Hiai2020}.
The map on cyclic orbits
\[
 x\xi\longmapsto xa^{1/2}\in L^2(M)
\]
is isometric, because both squared norms equal the functional evaluated at
$x^*x$.  Its extension intertwines the representations and has a reducing
range.  Decomposing into orthogonal cyclic reducing subspaces, of which
there are at most countably many, embeds the given representation into
$\widetilde K=L^2(M)^{\oplus\N}$.  These are standard consequences of the normal
representation theory of finite von Neumann algebras; see
\citet{Takesaki1979}.

Extend $Y$ by zero off the embedded reducing subspace.  The scalar
commutator bound is unchanged.  On the countable amplification, let $p_n$
be the projection onto the first $n$ copies.  Since $p_n$ commutes with the
common left action,
\[
 [p_nYp_n,L_x]=p_n[Y,L_x]p_n\qquad(x\in M).
\]
Proposition~\ref{prop:finite-copies} applies to the compression of $Y$ to
$p_n\widetilde K$, with constant $3c_1+2$.  For any fixed finite matrix test, the
corresponding compressed commutators converge strongly to the full
commutator as $n\to\infty$.  A uniform operator-norm bound passes to this
strong limit.  Restricting back to the original reducing subspace proves
the assertion for type $\mathrm{II}_1$ factors.

If $P$ is properly infinite, choose isometries $v_1,\ldots,v_h\in P$ with
orthogonal ranges.  The row $V=[v_1\ \cdots\ v_h]$ satisfies $V^*V=I_{K^h}$.
Apply the row and column derivation estimates of
Section~\ref{sec:preliminaries} to $x\mapsto[Y,x]$.  With
$\delta=g_P(Y)$, they give
\[
 \norm{YV-VY^{(h)}}\le c_{\mathrm{row}}\delta,
 \qquad
 \norm{Y^{(h)}V^*-V^*Y}\le c_{\mathrm{row}}\delta.
\]
For $X\in M_h(P)$, the identity
\begin{align*}
 [Y,VXV^*]
  &=(YV-VY^{(h)})XV^*
    +V[Y^{(h)},X]V^*\\
  &\hspace{2em}+VX(Y^{(h)}V^*-V^*Y)
\end{align*}
and $\norm{VXV^*}=\norm{X}$ imply
\[
 \norm{[Y^{(h)},X]}
 \le(1+2c_{\mathrm{row}})\delta\norm{X}.
\]
Here conjugation by the isometry $V:K^h\to K$ preserves the norm of the
middle term, regardless of whether its range is all of $K$.

Finally, if $P$ is a finite type I factor, average $Y$ over its compact
unitary group to obtain $Y_0\in P'$ with $\norm{Y-Y_0}\le\delta$.
Then
\[
 [Y^{(h)},X]=[(Y-Y_0)^{(h)},X],
\]
which gives the constant $2$.  These three cases exhaust the factors and
prove the proposition.
\end{proof}

\subsection{Central decomposition}

\begin{proposition}\label{prop:separable-estimate}
Let $P\subset\B(K)$ be any unital von Neumann algebra on a separable
Hilbert space.  Then
\[
 \norm{[Y^{(h)},X]}\le C\,g_P(Y)\,\norm{X}
 \qquad(Y\in\B(K),\ X\in M_h(P),\ h\ge1).
\]
\end{proposition}

\begin{proof}
Write $\delta=g_P(Y)$.  Average $Y$ over the compact unitary groups of the
finite-dimensional algebras generated by finite families of central
projections.  Direct these families by inclusion and take a cofinal weak
operator cluster limit.  The resulting operator $Y_1$ commutes with every
central projection, hence with $Z(P)$, and satisfies
\begin{equation}\label{eq:center-average}
 \norm{Y_1-Y}\le\delta,
 \qquad g_P(Y_1)\le3\delta.
\end{equation}
For the first inequality, each unitary being averaged belongs to $P$, so
its conjugation changes $Y$ by at most $\delta$.  The inequality is
preserved by averages and weak operator limits.  The second follows from
$g_P(Y_1)\le g_P(Y)+2\norm{Y_1-Y}$.

Use the central direct-integral decomposition on a separable Hilbert
space \citep{Takesaki1979}; precise section locators for the
factor fibers, decomposable operators, and norm formula are
\citet[Sections~III.1.6.1--4 and III.5.1.10--15]{Blackadar2006}:

\[
 K=\int^{\oplus}K_\alpha\,d\mu(\alpha),
 \qquad
 P=\int^{\oplus}P_\alpha\,d\mu(\alpha),
\]
where the $P_\alpha$ are factors almost everywhere and the fibers are
separable.  The center acts as the scalar diagonal algebra, so its
commutant consists of the decomposable operators.  In particular, $Y_1$
has fibers $Y_1(\alpha)$.  We verify carefully that
\begin{equation}\label{eq:fiber-commutator}
 g_{P_\alpha}\bigl(Y_1(\alpha)\bigr)\le3\delta
 \quad\text{for almost every }\alpha.
\end{equation}

The separable direct-integral decomposition admits countably many bounded
measurable sections whose values generate $P_\alpha$ almost everywhere;
fix such a family. Take all rational complex $*$-polynomials in
these sections, including the rational scalar constants.  For each such
section $b\in P$, form the contraction
\[
 c_b=b(1+b^*b)^{-1/2}\in P.
\]
On every fiber where the generators have the stated property, these
contractions are norm dense in the contraction ball of the generated
$C^*$-algebra.  To see this, the polynomials are norm dense, the displayed
transformation is norm continuous, and its range contains the open unit
ball: its inverse there is
\[
 c\longmapsto c(1-c^*c)^{-1/2}.
\]
The closed unit ball is the norm closure of the open one.  Kaplansky
density then makes these countably many fiber contractions strongly
dense in the unit ball of $P_\alpha$.

For each $c_b$, the bound $\norm{[Y_1,c_b]}\le3\delta$ implies the same
fiber bound outside a null set, by the essential-supremum norm formula.
Removing the union of these countably many null sets leaves a single
conull set on which every test holds.  Strong approximation by the
contraction tests passes the bound to every element of the fiber unit
ball, proving \eqref{eq:fiber-commutator}.

Apply Proposition~\ref{prop:factor-representations} on that conull set.
For any fixed $X\in M_h(P)$, decomposability of its entries and the same
norm formula give
\[
 \norm{[Y_1^{(h)},X]}
 \le3C_{\mathrm{fac}}\delta\norm{X}.
\]
Restoring $Y$ costs at most $2\delta\norm{X}$ by
\eqref{eq:center-average}, giving $C=3C_{\mathrm{fac}}+2$.
\end{proof}

\subsection{Removing separability}

\begin{proof}[Proof of Theorem~\ref{thm:universal}]
Let $P\subset\B(K)$ be arbitrary, and fix $h\ge1$, $X\in M_h(P)$, and a
unit vector $\xi=(\xi_1,\ldots,\xi_h)\in K^h$.  Let $A_X\subset P$ be
the unital $C^*$-algebra generated by the entries of $X$.  The
$C^*$-algebra in $\B(K)$ generated by these entries and $Y$ is norm
separable.  The closed span of its orbits of the finitely many $\xi_i$
is therefore a separable common reducing subspace $K_0$ for $A_X$ and
$Y$.  Denote restriction of $A_X$ to $K_0$ by $\sigma$, put
$Y_0=Y|_{K_0}$, and set $P_0=\sigma(A_X)''$.

For $a\in A_X$,
\[
 \norm{[Y_0,\sigma(a)]}\le g_P(Y)\norm{a}.
\]
The $C^*$-quotient norm of the image gives, for every $b\in\sigma(A_X)$,
\[
 \norm{[Y_0,b]}
 \le g_P(Y)\inf_{\sigma(a)=b}\norm{a}
 =g_P(Y)\norm{b}.
\]
By Kaplansky density the same scalar commutator bound holds on $P_0$.
Proposition~\ref{prop:separable-estimate}, applied to $P_0$ and
$\sigma_h(X)$, now yields
\[
 \norm{[Y_0^{(h)},\sigma_h(X)]}
 \le C g_P(Y)\norm{\sigma_h(X)}
 \le C g_P(Y)\norm{X}.
\]
Since $K_0$ reduces all the operators in this test and contains every
component of $\xi$, the restricted commutator on $\xi$ equals the original
one.  Taking the supremum over unit test vectors proves
\eqref{eq:universal}.  If $K=\{0\}$, the estimate is vacuous.
\end{proof}

\subsection{Derivations and similarity}

Theorem~\ref{thm:universal} now controls inner commutators in every
representation. To reach Kirchberg's criterion we must implement an
arbitrary bounded derivation. Cyclic-domain implementation supplies
finitely many columns at a time; the uniform commutator estimate lets
us pass from these finite collections to the full derivation.

\begin{proposition}\label{prop:all-derivations}
Let $A$ be a unital $C^*$-algebra and $\sigma:A\to\B(K)$ a unital
$*$-representation.  Every bounded $\sigma$-derivation
$\Delta:A\to\B(K)$ satisfies
\[
 \norm{\Delta}_{\mathrm{cb}}\le C\norm{\Delta}
\]
and is implemented by a bounded operator on $K$.
\end{proposition}

\begin{proof}
Decompose $K$ into cyclic reducing subspaces for $\sigma$:
\[
 K=\bigoplus_{i\in I}K_i.
\]
Such a decomposition is obtained by taking a maximal orthogonal family
of cyclic reducing subspaces; a nonzero vector in its orthogonal
complement would generate another one.  The index set $I$ can have any
cardinality.  Let $q_i$ denote the projection onto $K_i$.

The map $a\mapsto\Delta(a)|_{K_i}$ is a rectangular derivation with cyclic
domain representation $\sigma|_{K_i}$ and full output representation
$\sigma$.  Lemma~\ref{lem:cyclic-implementation} gives a bounded map
$V_i:K_i\to K$ such that
\[
 \Delta(a)|_{K_i}=V_i\sigma(a)|_{K_i}-\sigma(a)V_i.
\]
For a finite subset $F\subset I$, put
\[
 q_F=\sum_{i\in F}q_i,
 \qquad V_F=\sum_{i\in F}V_iq_i\in\B(K).
\]
Every $V_F$ is bounded, and the reducing property gives the exact identity
\begin{equation}\label{eq:finite-column-identity}
 [V_F,\sigma(a)]=\Delta(a)q_F\qquad(a\in A).
\end{equation}
Only the scalar commutator norm, rather than the norm of $V_F$, needs to
be uniform in $F$.  Indeed, taking the infimum over lifts of an element
$b\in\sigma(A)$ in \eqref{eq:finite-column-identity} gives
\[
 \norm{[V_F,b]}\le\norm{\Delta}\norm{b}.
\]
This uses the $C^*$-quotient norm of the $*$-image and does not require
$\sigma$ to be faithful.  Kaplansky density then implies
\[
 g_{\sigma(A)''}(V_F)\le\norm{\Delta}.
\]

Apply Theorem~\ref{thm:universal} to $V_F$ and $\sigma(A)''$.  For every
$X\in M_h(A)$,
\begin{align*}
 \norm{\Delta_h(X)q_F^{(h)}}
   &=\norm{[V_F^{(h)},\sigma_h(X)]}\\
   &\le C\norm{\Delta}\norm{\sigma_h(X)}
    \le C\norm{\Delta}\norm{X}.
\end{align*}
The projections $q_F$ converge strongly to the identity as a net indexed
by finite subsets of $I$.  Hence the displayed bound passes to
$\Delta_h(X)$ for every fixed $h$ and $X$.  This proves the claimed
complete boundedness.  Lemma~\ref{lem:cb-implementation} then supplies a
bounded implementing operator.
\end{proof}

\begin{proof}[Proof of Theorem~\ref{thm:similarity}]
For every faithful unital $*$-representation of $A$, each bounded
derivation into the ambient bounded operators is inner by
Proposition~\ref{prop:all-derivations}.  Thus condition~(2) of
Theorem~\ref{thm:kirchberg} holds for $A$.  Its equivalent similarity
property gives, for every bounded unital homomorphism
$\pi:A\to\B(H)$, a bounded invertible $S$ with bounded inverse such
that $S\pi(\cdot)S^{-1}$ is a $*$-homomorphism.
The identity isomorphism treats the zero Hilbert space as well.
\end{proof}

\subsection{Hyperreflexivity}\label{sec:hyperreflexivity}

The uniform commutator estimate also controls distance to any von Neumann
algebra. We apply it to the commutant, whose projections describe the
algebra's invariant subspaces.

\begin{proof}[Proof of Corollary~\ref{cor:hyperreflexivity}]
Put \(P=M'\). Since \(M\) is self-adjoint, its invariant subspaces are
reducing and their orthogonal projections are exactly the projections in
\(P\). Hence \(\alpha_M\) is the usual off-diagonal distance seminorm.
For every projection \(e\in P\), the two off-diagonal blocks give
\[
 \|[T,e]\|=\max\{\|(1-e)Te\|,\|eT(1-e)\|\}\le\alpha_M(T),
\]
where the last inequality also uses the projection \(1-e\).
Spectral step approximations of a self-adjoint contraction \(a\in P\)
express it, modulo a scalar, as a norm limit of linear combinations of
projections in \(P\) whose coefficients have total absolute value at most two.
Thus \(\|[T,a]\|\le2\alpha_M(T)\). The real and imaginary parts of
any contraction in \(P\) are self-adjoint contractions, so
\(g_P(T)\le4\alpha_M(T)\).

Arveson's commutant-distance formula, in the form recalled in
\citet[Equation~(2.1)]{ChristensenSinclairSmithWhite2010}, gives
\[
 2\dist(T,M)=\|\operatorname{ad}_T|_P\|_{\mathrm{cb}},
 \qquad \operatorname{ad}_T(a)=[T,a],
\]
since \(P'=M\). Taking the supremum over finite matrix sizes in
Theorem~\ref{thm:universal} therefore yields
\[
 2\dist(T,M)\le Cg_P(T)\le4C\alpha_M(T),
\]
which proves Equation~\eqref{eq:universal-hyperreflexivity}.
\end{proof}
